\section{Disk-Backed Growable Vectors with Stable Virtual Addresses}
\label{sec:disk_backed_growable_vectors}

High-resolution Mandelbrot rendering, perturbation theory (\cref{sec:perturbation-concept}), and reference-orbit computation (\cref{sec:ref-orbit-calc}) impose
extreme demands on memory management. At deep zoom levels, a single authoritative
reference orbit may contain tens of billions of elements; for example, an orbit
with $2.7\times10^{10}$ iterations already implies a logical memory footprint far
beyond what can be held resident in ordinary RAM configurations. At 16 bytes per
complex-double record, $2.7\times10^{10}$ records require $432$ GB, approximately
$402$ GiB, before any coefficient tables or allocator overhead.
At this scale, both the cost of allocation and the cost of relocation become
dominant design constraints.

Several subsystems in \FractalShark{}—including orbit logging (see the compressed orbit format in \cref{subsec:ref-orbit-compression-reuse}), metadata streams,
debug outputs, and the custom heap allocator—require \emph{growable} storage
while simultaneously maintaining \emph{stable virtual addresses} for previously
returned pointers. This requirement arises naturally when the program hands out
interior pointers into a buffer (e.g., into an orbit record array or a heap
arena). Any reallocation-and-copy strategy would invalidate those pointers and
would require copying large amounts of data, which can dominate the calculation at multi-billion-element scales.

In addition to scale and pointer stability, persistence is a first-class
concern. Reference orbits are expensive to compute (\cref{subsec:ref-orbit-gpu}) but comparatively cheap to
reuse. If orbit data can be left resident on disk and later reloaded simply by
mapping the same backing file back into the process address space, recomputation
can be avoided entirely. This persistence enables efficient multi-pass rendering workflows,
interactive debugging, and reuse of authoritative reference orbits across program
invocations.

A further motivation is memory pressure.  Reference orbits, LA~v2
coefficient tables, and compressed-orbit metadata can collectively
grow to tens of gigabytes at extreme zoom depths.  Backing these
structures with on-disk files rather than anonymous (pagefile-backed)
memory avoids exhausting the system-wide virtual-memory commit charge,
which on Windows is limited by RAM plus pagefile size.  With
file-backed storage, the operating system demand-pages data into
physical memory as needed and evicts it back to the file under
pressure. A small active working set can therefore use much less resident RAM
than the logical data set, although a workload touching most records still incurs
memory pressure and disk I/O.

To address these requirements, \FractalShark{} provides a
\code{GrowableVector<EltT>} abstraction that combines (i) a fixed,
virtually-contiguous address range reserved up front and (ii) incremental growth
within that range by making additional capacity accessible on demand. Windows uses
commit or section extension; Linux uses memory mappings and protection changes.
Growth is performed
in place, so the base address never changes and no copying is required as
capacity increases, up to the chosen view size. The stable address applies during
the live mapping, rather than guaranteeing the same address in a later process.

\paragraph{Worked memory accounting example.}
Suppose the vector reserves a 16 GiB virtual view, enables 256 MiB of capacity,
contains 10 MiB of records, and has recently touched 4 MiB of pages. These are four
different quantities. Growing capacity to 512 MiB leaves pointers into the retained
records unchanged; it does not mean that all 512 MiB immediately reside in RAM.
Linux temporary backing can have a 16 GiB sparse logical file length while consuming
few disk blocks initially. Actual residence and disk allocation depend on access
and operating-system paging.

The abstraction supports two complementary backing modes. The following APIs
describe the Windows implementation; the Linux mapping is explained at the end:

\begin{enumerate}
  \item \textbf{Anonymous (pagefile-backed) growth.} A large virtual region is
  reserved once via \code{VirtualAlloc(..., MEM\_RESERVE, ...)} and pages are
  subsequently committed within that region using
  \code{VirtualAlloc(base, bytes, MEM\_COMMIT, ...)}. Because each commit
  targets the original base address, successful growth preserves
  \code{m\_Data} and does not require relocation.

  \item \textbf{File-backed (disk-backed) growth.} The reserved address range is
  backed by a disk file and mapped into the process using a section object. In
  this mode, pages are demand-paged from the backing file and become resident
  only as accessed, substantially reducing peak commit requirements when the
  logical capacity greatly exceeds physical memory. When persistence is enabled,
  the same backing file can later be remapped to recover the stored data without
  recomputation, directly supporting the orbit-reuse strategy described in \cref{subsec:ref-orbit-reuse}.
\end{enumerate}

\subsection{Stable-address growth via reserve-and-commit}
\label{sec:reserve_commit}

\begin{figure}[!htbp]
\centering
\begin{tikzpicture}[
    >=Stealth,
    cell/.style={draw, minimum width=0.55cm, minimum height=0.7cm, font=\tiny},
    committed/.style={cell, fill=blue!20},
    reserved/.style={cell, fill=gray!10, pattern=north east lines, pattern color=gray!30},
    label/.style={font=\scriptsize},
]
  % Reserved region
  \foreach \i in {0,...,19} {
    \ifnum\i<12
      \node[committed] (c\i) at (\i*0.6, 0) {};
    \else
      \node[reserved] (c\i) at (\i*0.6, 0) {};
    \fi
  }

  % Base pointer
  \draw[<-, thick] (c0.south) -- ++(0, -0.6) node[below, label] {\code{m\_Data} (stable)};

  % Capacity marker
  \draw[<->, thick, blue!70!black] (c0.north west) ++(0,0.15) -- ++(7.2, 0)
    node[midway, above, label, blue!70!black] {Accessible capacity (not used size)};

  % Reserved marker
  \draw[<->, thick, gray] (c12.north west) ++(0,0.15) -- ++(4.8, 0)
    node[midway, above, label, gray] {Reserved (uncommitted)};

  % Growth arrow
  \draw[->, thick, orange, dashed] (c11.east) ++(0.1, -0.75) -- ++(1.4, 0)
    node[right, label, orange] {Enable more capacity};

  % Total region bracket
  \draw[decorate, decoration={brace, amplitude=5pt, mirror}, thick]
    (c0.south west) ++(0, -1.2) -- ++(12.0, 0)
    node[midway, below=5pt, label] {Virtual address reservation (no reallocation)};
\end{tikzpicture}
\caption[\code{GrowableVector} reserve-and-commit memory layout]{\code{GrowableVector} reserve-and-commit memory layout.
  A large virtual region is reserved at initialization, establishing a
  stable base pointer \code{m\_Data}. Accessible capacity grows incrementally
  (blue, left), within the reserved region (hatched, right). Capacity, used
  elements, and resident pages are distinct: making a page accessible does not
  require it to remain resident. Pointers into retained records remain valid
  across successful growth operations.}
\label{fig:growable-vector}
\end{figure}

At initialization time, the vector establishes an address anchor:
\begin{itemize}
  \item In anonymous mode, the anchor is created by reserving a large region
  (typically sized to a fraction of physical memory, or an explicit override)
  and recording the returned base pointer \code{m\_Data}.
  \item In file-backed mode, the anchor is created by mapping a section at a
  base address and thereafter requiring all remaps/extends to preserve that base.
\end{itemize}

Growth proceeds by increasing the logical capacity \code{m\_CapacityInElts}
without moving data. In anonymous mode, this corresponds to committing
additional pages within the reserved region on Windows, or changing page
protections on Linux. Windows file-backed growth extends the section; Linux
persistent growth extends the file within a fixed mapped view, while temporary
sparse backing already has the full view's logical length. In both cases the implementation \emph{never} allocates a second buffer and
copy: the address stability is enforced by explicitly requesting the existing
base address and treating any base-address change as a hard error.

\subsection{File-backed mapping for low-commit large buffers}
\label{sec:file_backed_mapping}

When a filename is provided (or a temporary filename is generated), the Windows
vector opens or creates a backing file and constructs a section object over it. It uses
native NT section APIs (e.g., \code{NtCreateSection},
\code{NtMapViewOfSection}, and \code{NtExtendSection}) to obtain a
reserve-style mapping suitable for later extension.

Conceptually, the pipeline is:
\begin{enumerate}
  \item \textbf{Open/Create backing file.} The file can be persistent (saved
  results) or temporary (e.g., created with delete-on-close semantics when the
  data is not meant to be retained).
  \item \textbf{Create section.} A section is created with read/write
  protection. For growable vectors, the section is configured so that the view
  can be reserved and later extended.
  \item \textbf{Map a view to establish the base address.} The first mapping
  produces the stable anchor pointer \code{m\_Data}. Subsequent growth must
  preserve this base.
  \item \textbf{Extend on demand.} When capacity increases, the section is
  extended to the new byte size. Because the base mapping remains fixed, all
  previously returned pointers remain valid.
\end{enumerate}

This approach has two key practical benefits:
\begin{itemize}
  \item \textbf{Stable pointers with large logical capacity.} The program may
  reserve and grow to very large capacities (hundreds of GiB or more) without
  the pointer invalidation inherent in \code{realloc}-style strategies.
  \item \textbf{Reduced peak commit pressure.} The OS can keep most of the
  logical buffer non-resident until touched; pages are faulted in as accessed
  and can be evicted back to disk under memory pressure. This behavior is particularly
  effective when the workload has a large addressable dataset but a much smaller
  active working set.
\end{itemize}

\subsection{Trimming and persistence policies}
\label{sec:trim_persist}

The vector supports multiple policies that control whether data is persisted:
\begin{itemize}
  \item \textbf{Persistent save (\code{EnableWithSave}).} Upon \code{Trim()},
  the backing file is truncated to match the used size, ensuring the on-disk
  representation reflects the logical content exactly.
  \item \textbf{Temporary file-backed (\code{EnableWithoutSave}).} Upon
  \code{Trim()}, the view can be remapped to match the used size while
  \emph{reusing the same base address}. This remapping reduces the mapped span without
  sacrificing the addresses of retained records. Windows deletes the temporary
  file on close; Linux unlinks it after opening, and its storage survives through
  the live descriptor and mapping.
  \item \textbf{Anonymous (\code{DontSave}).} Trimming is a logical operation
  (used size changes) while the reserved region remains as the stable anchor.
\end{itemize}

\subsection{Using the same mechanism for the heap arena}
\label{sec:heap_uses_growable_vector}

The custom heap allocator builds its arena on top of the same
\code{GrowableVector<uint8\_t>} abstraction. The heap is initialized by
creating a temporary file-backed growable vector for the heap arena and making
an initial chunk accessible. When the allocator needs more space, it expands the
arena by growing the underlying vector in-place and then updating the heap's
wilderness block (coalescing if possible). Because the arena grows within a
reserved and stable virtual range, all heap pointers returned to the program
remain valid across heap expansions, and expansion does not require copying or
relocating the heap.

\subsection{Portability and backing-file considerations}
\label{sec:disk_backed_portability_notes}

Windows and Linux implement the same stable-address contract with different
operating-system mechanisms. On Windows, anonymous storage uses
\code{VirtualAlloc} reserve/commit, while file-backed storage uses native section
objects and explicitly extends ordinary backing files. Reserving an address range,
charging committed storage, and making a page resident are separate events.

On Linux, anonymous storage reserves an \code{mmap} region with
\code{PROT\_NONE}; growth enables pages with \code{mprotect}. File-backed storage
uses a fixed \code{MAP\_SHARED} view. Persistent files grow through
\code{ftruncate}. Temporary mode gives a sparse file the full view's logical
length once, unlinks the file after opening, and subsequently grows capacity within
that existing view. A sparse file allocates disk blocks for written regions rather
than requiring physical storage for its complete logical length.

Temporary Linux trimming uses \code{mremap} without permission to move the mapping
and checks that the address stayed fixed. Persistent trimming truncates the backing
file to the used byte count. Neither platform promises unlimited growth: capacity
must fit the selected view, and shrinking removes storage beyond the retained
records. Linux protection changes also do not have the same commit-accounting
semantics as Windows commits.

The common ownership rule is simpler than either implementation: existing live
records keep their addresses during successful growth, and callers release those
records before destroying their backing vector. Persistence retains data, not a
particular virtual address across application runs.

The preceding chapters have covered the full rendering pipeline from iteration
kernels through memory management.  The next chapter (\cref{sec:future-directions}) surveys alternative and
complementary techniques that informed \FractalShark{}'s design choices and
that represent promising directions for future work.
